\cleardoublepage
\phantomsection
\pdfbookmark{Sommario}{Sommario}
\begingroup
\let\clearpage\relax
\let\cleardoublepage\relax
\let\cleardoublepage\relax

\chapter*{Sommario}
Il presente elaborato documenta lo stage di 320 ore svolto dalla studentessa Giulia Barzon presso Wintech S.p.A., a conclusione del Corso di Laurea in Informatica dell'Università degli Studi di Padova.

L'obiettivo principale dello stage è stato l'analisi, la valutazione e l'integrazione di una piattaforma di \textit{Data Catalog} all'interno dell'infrastruttura aziendale, 
per mantenere un catalogo dati centralizzato e garantire il rispetto delle normative sulla privacy, come il GDPR, e delle policy aziendali.\\

Partendo dall'analisi dei requisiti aziendali e da uno studio comparativo delle principali piattaforme di \textit{Data Catalog} 
presenti sul mercato, è stata individuata la piattaforma open-source OpenMetadata come lo strumento più idoneo a soddisfare le esigenze di Wintech.

La fase successiva si è concentrata sul deployment e sulla configurazione del software, testandone in modo completo le funzionalità chiave: la \textit{Data Ingestion} 
dalle diverse fonti aziendali, la \textit{Data Discovery} e la \textit{Data Quality}. 
Particolare attenzione è stata rivolta alla sicurezza e alla \textit{Data Governance}, validando l'efficacia del motore di individuazione e classificazione automatica dei dati sensibili e personali.\\

Per completare il lavoro svolto, parte delle ore dello stage è stata dedicata allo studio di Microsoft Power BI e dei principi della \textit{Business Intelligence}, 
dimostrando come l'integrazione tra strumenti di BI e piattaforme di \textit{Data Catalog} possa migliorare la gestione dei dati aziendali, garantendo la loro qualità, sicurezza e tracciabilità.

%\vfill

%\selectlanguage{english}
%\pdfbookmark{Abstract}{Abstract}
%\chapter*{Abstract}

%\selectlanguage{italian}

\endgroup

\vfill
